\documentclass[12pt,a4paper]{ctexart}
\usepackage[top=2.5cm,bottom=2.5cm,left=2.5cm,right=2.5cm]{geometry}
\usepackage{amsmath,amssymb,amsthm,physics}
\usepackage{graphicx,float,booktabs,array,caption}
\usepackage{hyperref,xcolor,enumitem}
\usepackage[numbers,sort&compress]{natbib}
\hypersetup{colorlinks=true,linkcolor=blue,citecolor=blue,urlcolor=blue}
\theoremstyle{definition}
\newtheorem{definition}{定义}[section]
\newtheorem{remark}{注记}[section]
\newtheorem{theorem}{定理}[section]

\title{\heiti 格点QCD中的OPE算符：\\
从短距展开到光锥因子化的系统框架}
\author{基于本库文献的综合解析}
\date{\today}

\begin{document}
\maketitle

\begin{abstract}
\noindent
算符乘积展开（Operator Product Expansion, OPE）是QCD中将非局域算符的乘积在短距离极限下展开为局域算符之和的核心理论工具。自Wilson于1969年提出这一框架以来，OPE已经成为连接微扰QCD的短距离物理与非微扰强子矩阵元的理论基石。在格点QCD的部分子物理计算中，OPE扮演着多层面的关键角色：（1）标准短距离OPE——将两个接近的场算符展开为局域twist-2算符的级数，其矩阵元给出PDF的Mellin矩；（2）光锥OPE（light-ray OPE）——采用不同的幂次计数方案，将涉及大动量标度的非局域算符重求和为``光射线''结构，这是LaMET和赝PDF因子化定理的理论基础；（3）重正子分析——通过OPE中Wilson系数的阶乘发散揭示twist-3（领先幂次）修正的函数形式；（4）靶质量修正——从OPE的运动学幂次展开中导出有限强子质量的$M^2/(P^z)^2$修正。本文系统梳理这四重OPE框架：从Wilson的原始短距离OPE出发，经由Izubuchi等人的因子化定理OPE证明，到Braun-Vladimirov-Zhang的光线OPE和重正子分析，最终展示OPE在预测准PDF和赝PDF中幂次修正的$x$依赖函数形式方面的核心作用。
\end{abstract}

\tableofcontents
\newpage

\section{引言：OPE——QCD的微扰-非微扰桥梁}

算符乘积展开（OPE）是Wilson于1969年提出的QCD核心理论工具。其基本思想是：两个（或多个）场算符的乘积在\textbf{短距离}极限下可以展开为一系列\textbf{局域算符}的和，每个局域算符乘以一个包含短距离奇点的\textbf{Wilson系数函数}（coefficient function）\cite{JaffeManohar1990,Izubuchi2018,Braun2019power}。

在数学上，OPE的形式为\cite{Izubuchi2018}：

\begin{definition}[标准短距离OPE]
\begin{equation}
\boxed{\mathcal{O}_1(z) \mathcal{O}_2(0) \xrightarrow{z \to 0} \sum_n C_n(z, \mu) \, \mathcal{O}_n^{\text{local}}(0)},
\label{eq:standard_OPE}
\end{equation}
其中$C_n(z, \mu)$是Wilson系数（微扰可计算，在$z \to 0$时具有奇点），$\mathcal{O}_n^{\text{local}}$是量纲递增的局域算符，$\mu$是因子化标度。
\end{definition}

\textbf{量纲分析决定了展开的收敛速度：}算符$\mathcal{O}_n^{\text{local}}$的量纲越高，其在展开中的贡献被$(z\Lambda_{\text{QCD}})^{d_n-d_1}$的幂次因子压低。因此，最低量纲（``leading-twist''）算符的贡献主导了$z \to 0$极限下的行为。

\textbf{在格点QCD中，OPE扮演着三重独特且互补的角色}，这正是本文的论述主线：

\begin{enumerate}
    \item \textbf{传统角色（PDF矩的计算）：}将非局域光锥算符展开为局域twist-2算符，提取Mellin矩$a_n(\mu) = \int dx\, x^{n-1} q(x, \mu)$\cite{Ji2013Euclidean}。由于格点上高维算符的噪声和混合问题，实际只能可靠计算前两三个矩。

    \item \textbf{现代角色（LaMET因子化）：}Izubuchi等人\cite{Izubuchi2018}利用OPE严格证明了准PDF与光锥PDF之间的因子化定理——这是LaMET理论合法性的基石。OPE将准PDF中的``硬''（短距离）模式和``软''（长距离）模式系统地分离。

    \item \textbf{精度角色（幂次修正）：}Braun、Vladimirov和Zhang\cite{Braun2019power}利用\textbf{光线OPE}（light-ray OPE）和重正子分析，预言了准PDF和赝PDF中领先幂次（twist-3）修正对Bjorken-$x$的完整函数依赖性。
\end{enumerate}

\section{标准短距离OPE与PDF的Mellin矩}

\subsection{从光锥关联到局域算符的Taylor展开}

在LaMET方法出现之前，格点QCD计算PDF的唯一途径是通过OPE将光锥关联展开为局域算符，然后提取其矩阵元（Mellin矩）。具体而言，将光锥PDF定义中的非局域算符在$\xi^- \to 0$处展开\cite{Ji2013Euclidean}：

\begin{equation}
\begin{aligned}
\bar{\psi}(\xi^-)\gamma^+ \mathcal{U}(\xi^-,0)\psi(0) &= \sum_{n=0}^{\infty} \frac{(\xi^-)^n}{n!} \bar{\psi}(0)\gamma^+ (iD^+)^n \psi(0) \\
&= \sum_{n=0}^{\infty} \frac{(\xi^-)^n}{n!} \mathcal{O}^{+\ldots+}(0)。
\end{aligned}
\label{eq:Taylor_OPE}
\end{equation}

其中$\mathcal{O}^{\mu_1\ldots\mu_n} = \bar{\psi}\gamma^{(\mu_1}iD^{\mu_2}\ldots iD^{\mu_n)}\psi - \text{trace}$是\textbf{无迹对称化的twist-2局域算符}。取核子矩阵元并作Fourier变换，得到PDF与局域矩阵元之间的关系：

\begin{equation}
\int_0^1 dx\, x^{n-1} q(x, \mu^2) = a_n(\mu^2), \qquad \langle P|\mathcal{O}^{\mu_1\ldots\mu_n}(\mu^2)|P\rangle = 2a_n(\mu^2)(P^{\mu_1}\ldots P^{\mu_n} - \text{trace})。
\label{eq:moment_relation}
\end{equation}

\subsection{格点上Mellin矩计算的困难}

\textbf{为什么格点上只能计算前两三个矩？}Ji在LaMET的原始论文中总结了其中的技术原因\cite{Ji2013Euclidean}：

\begin{enumerate}
    \item \textbf{算符混合：}格点的超立方对称群（代替连续Lorentz群$O(4)$）具有较少的不可约表示。因此，在连续理论中属于不同自旋表示的局域算符在格点上会发生混合。高自旋（高$n$）算符的混合模式极为复杂。

    \item \textbf{幂次发散的混合：}在格点截断方案下，高维twist-2算符可能与低维算符发生\textbf{幂次发散}的混合——即混合系数$\sim 1/a^{\Delta d}$，其中$\Delta d$是量纲差。这使得在未重整化的情况下，高矩的提取在数值上不可能。

    \item \textbf{信噪比随$n$急剧恶化：}高维算符涉及更多协变导数$D^\mu$，在Monte Carlo模拟中产生极大的统计涨落。
\end{enumerate}

\textbf{LaMET如何绕过这些困难？}Ji指出\cite{Ji2013Euclidean,Ji2017More}：通过保持算符的\textbf{非局域形式}——而非在计算前将其展开为局域算符的级数——可以完全避免算符混合和幂次发散的问题。非局域算符的UV行为由Wilson线的自能（线性发散$\sim e^{-\delta m|z|}$）和算符顶点的对数发散控制——这些发散是\textbf{乘性的}且已被充分理解。

\subsection{$g_1$求和规则的OPE分析：一个历史范例}

Jaffe和Manohar在其1990年的经典论文\cite{JaffeManohar1990}中对$g_1$结构函数的OPE分析构成了理解这一框架的最佳范例。$g_1$求和规则的OPE分析``惊人地简单''\cite{JaffeManohar1990}：

\begin{quote}
``The OPE analysis of $\int dx\, g_1(x,Q^2)$ in QCD is in one sense remarkably simple: this sum rule projects out a unique operator, $\bar{\psi}\gamma^\mu\gamma_5\mathcal{Q}^2\psi$, where $\mathcal{Q}$ is the quark charge matrix; there is no gauge invariant gluonic operator with this dimension and spin and therefore no operator mixing problem.''
\end{quote}

\textbf{关键的物理洞察：}非单态（同位旋矢量）轴矢流具有\textbf{零反常量纲}——在微扰QCD中不演化。单态（味单态）轴矢流具有从双圈图级别开始的非零反常量纲（涉及图2中的胶子中间态），导致极化的缓慢``转移''——夸克自旋通过反常过程缓慢转移到夸克轨道角动量\cite{JaffeManohar1990}：

\begin{equation}
\Delta\Sigma(Q^2) = \exp\left[ -\frac{2}{\pi}\frac{\alpha_s(Q^2) - \alpha_s(Q_0^2)}{33 - 2N_f} \right] \Delta\Sigma(Q_0^2)。
\label{eq:singlet_evolution}
\end{equation}

\section{光线OPE（Light-Ray OPE）：非局域算符的``重求和''展开}

\subsection{短距离OPE的局限与光线OPE的提出}

标准短距离OPE要求$|z| \to 0$（以$\Lambda_{\text{QCD}}^{-1}$为单位），这使得它在格点上的应用被限制在非常小的$z$区域。然而，LaMET和赝PDF方法中使用的非局域算符涉及有限甚至较大的空间分离$z$，此时标准OPE的收敛性存疑。

\textbf{光线OPE（light-ray OPE）}\cite{Braun2019power}通过采用\textbf{不同的幂次计数方案}解决了这个问题：

\begin{definition}[光线OPE的幂次计数]\cite{Braun2019power}
\begin{equation}
\boxed{z = \mathcal{O}(\eta), \qquad p = \mathcal{O}(\eta^{-1}), \qquad \eta \to 0}。
\label{eq:light_ray_power_counting}
\end{equation}
在这种计数方案下，$zP^z = \mathcal{O}(1)$（即Ioffe时间保持有限），但$z\Lambda_{\text{QCD}} \to 0$。leading-twist算符的矩阵元标度为$\langle\mathcal{O}_n\rangle \sim \eta^{-n}\Lambda_{\text{QCD}}^n$，使得级数中的每一项$z^n \langle\mathcal{O}_n\rangle \sim \mathcal{O}(1)$，必须被\textbf{重求和到所有阶}。
\end{definition}

\textbf{光线OPE的显式形式为}\cite{Braun2019power}：

\begin{equation}
\boxed{\bar{q}(zv)\slashed{v}[zv,0]q(0) = \int_0^1 d\alpha\, H_\parallel(z, \alpha, \mu_F) \, \Pi^{\mu_F}_{\text{l.t.}}[\bar{q}(\alpha zv)\slashed{z}q(0)] + \ldots},
\label{eq:light_ray_OPE}
\end{equation}

其中$\Pi^{\mu_F}_{\text{l.t.}}[\ldots]$是\textbf{leading-twist投影算符}，定义为重正化局域leading-twist算符的生成函数：

\begin{definition}[Leading-Twist投影算符]\cite{Braun2019power}
\begin{equation}
\boxed{\Pi^{\mu_F}_{\text{l.t.}}[\bar{q}(zv)\slashed{v}q(0)] = \sum_{n=1}^{\infty} \frac{z^{n-1}}{(n-1)!} v_{\mu_1}\ldots v_{\mu_n} \mathcal{O}_n^{\mu_1\ldots\mu_n}(0)},
\label{eq:leading_twist_projection}
\end{equation}
其中$\mathcal{O}_n^{\mu_1\ldots\mu_n}(z) = \bar{q}(0)\gamma^{(\mu_1}\overleftarrow{D}^{\mu_2}\ldots\overleftarrow{D}^{\mu_n)}q(0)$是重正化的无迹对称化局域算符，其矩阵元给出Mellin矩：$\langle P|\mathcal{O}_n^{\mu_1\ldots\mu_n}|P\rangle = 2a_n(\mu)(P^{\mu_1}\ldots P^{\mu_n} - \text{trace})$。
\end{definition}

\textbf{光线OPE与标准OPE的本质区别：}

\begin{table}[H]
\centering
\caption{标准短距离OPE vs 光线OPE}
\label{tab:OPE_comparison}
\begin{tabular}{lcc}
\toprule
特性 & 标准OPE & 光线OPE \\
\midrule
小参量     & $z$（以$\Lambda_{\text{QCD}}^{-1}$为单位的空间分离） & $z$和$1/p$同步趋于零 \\
Ioffe时间  & $\nu = zP^z \to 0$ & $\nu = \mathcal{O}(1)$（保持有限） \\
展开形式  & 截断到有限阶 & \textbf{重求和到所有阶} \\
算符标度  & $\langle\mathcal{O}_n\rangle \sim \Lambda_{\text{QCD}}^n$ & $\langle\mathcal{O}_n\rangle \sim \eta^{-n}\Lambda_{\text{QCD}}^n$ \\
适用场景  & 局域算符、小$z$ & LaMET/赝PDF的非局域算符 \\
截断误差  & 高阶算符（量纲$>2+N$） & 高阶twist（twist $\ge 4$） \\
\bottomrule
\end{tabular}
\end{table}

\subsection{从光线OPE到因子化定理}

将光线OPE应用于核子矩阵元，直接导出LaMET的因子化定理\cite{Braun2019power}：

\begin{equation}
\begin{aligned}
I_\parallel(z^2 v^2, pvz, \mu \sim 1/|vz|) &= \int_0^1 d\alpha\, H_\parallel(z, \alpha, \mu_F) \left[1 - \frac{1}{4}\frac{m^2 v^2}{(pv)^2} z\frac{d}{dz} + \ldots\right] \\
&\quad \times \mathcal{I}(\alpha z pv, \mu_F) + \ldots,
\end{aligned}
\label{eq:factorization_from_OPE}
\end{equation}

其中$\mathcal{I}(\alpha z pv, \mu_F) = \int_{-1}^1 dx\, e^{ix\alpha zpv} q(x, \mu_F)$是光锥PDF的Fourier变换（Ioffe-时间分布，ITD），$H_\parallel$是微扰可计算的Wilson系数（$\delta(1-\alpha) + \mathcal{O}(\alpha_s)$），$m^2 v^2/(pv)^2$项是靶质量修正。

\section{OPE与LaMET因子化定理的严格证明}

\subsection{Izubuchi等人的OPE证明}

Izubuchi、Ji、Jin、Stewart和Zhao\cite{Izubuchi2018}利用OPE给出了\textbf{LaMET因子化定理的严格证明}。他们的核心论证分为以下几个步骤：

\begin{enumerate}
    \item \textbf{算符识别：}准PDF算符$\mathcal{O}_\Gamma(z) = \bar{\psi}(z)\Gamma U(z,0)\psi(0)$在短距离$z \to 0$下的OPE展开，其leading-twist贡献与光锥PDF算符具有相同的自旋结构。

    \item \textbf{因子化结构的推导：}通过OPE将准PDF矩阵元中的``硬''模式（$\sim 1/z$量级的动量）分离到Wilson系数$C$中，剩余的``软''模式由光锥PDF的非微扰矩阵元编码：

    \begin{equation}
    \tilde{h}(zP^z, z^2, a^{-1}) = \sum_n C_n(z^2, a^{-1}, \mu) \, \langle P|\mathcal{O}_n^{\text{local}}(\mu)|P\rangle + \mathcal{O}(z^2\Lambda_{\text{QCD}}^2)。
    \label{eq:OPE_factorization}
    \end{equation}

    \item \textbf{IR匹配：}准PDF和光锥PDF的OPE具有\textbf{相同的红外结构}（两者的非微扰矩阵元都由相同的局域算符$\mathcal{O}_n^{\text{local}}$的矩阵元给出）。因此，两者的差异完全由UV Wilson系数$C_n$的不同来描述——这些差异正是匹配核的来源。

    \item \textbf{幂次修正的分类：}高阶twist（twist-3、twist-4等）算符的贡献被$(z\Lambda_{\text{QCD}})^{t-2}$或等价的$(\Lambda_{\text{QCD}}/P^z)^{t-2}$因子压低。
\end{enumerate}

\subsection{准PDF与赝PDF的等价性：来自OPE的视角}

Izubuchi等人\cite{Izubuchi2018}的OPE分析同样证明了\textbf{准PDF和赝PDF方法的等价性}：

\begin{quote}
``There is a unique factorization formula that matches the quasi-PDF, spatial correlator and pseudo-PDF to the PDF. Since their factorizations into the PDF all require small distances and have large nucleon momentum, the setup for their lattice calculations must also be the same. Therefore, the LaMET and pseudo distribution approaches are in principle equivalent to each other.''
\end{quote}

从OPE的角度看，这两种方案仅是同一展开在\textbf{不同``空间''中执行}的结果：LaMET在动量空间中执行（Fourier变换后再匹配），赝PDF在坐标空间中执行。两者的Wilson系数通过Fourier变换相联系。

\subsection{约化Ioffe-时间分布的短距离行为}

Izubuchi等人\cite{Izubuchi2018}进一步利用OPE解释了\textbf{约化赝PDF在小的$z^2$处具有弱对数依赖性}这一数值观察。在短距离$z^2 \to 0$下：

\begin{equation}
\frac{\tilde{Q}(\zeta, \mu^2 z^2)}{\tilde{Q}(0, \mu^2 z^2)} = \sum_n \frac{C_n(\mu^2 z^2)}{C_0(\mu^2 z^2)} \frac{(-i\zeta)^n}{n!} a_{n+1}(\mu) + \mathcal{O}(z^2\Lambda_{\text{QCD}}^2)。
\label{eq:ratio_OPE}
\end{equation}

其中的NLO Wilson系数$C_n(\mu^2 z^2)$被显式计算为\cite{Izubuchi2018}：

\begin{equation}
C_n(\mu^2 z^2) = 1 + \frac{\alpha_s C_F}{2\pi}\left[ \left(\frac{3+2n}{2+3n+n^2} + 2H_n\right) \ln\frac{\mu^2 z^2 e^{2\gamma_E}}{4} + \frac{5+2n}{2+3n+n^2} + 2(1-H_n)H_n - 2H^{(2)}_n \right]。
\end{equation}

\textbf{关键的数值发现：}$\tilde{Q}(0, \mu^2 z^2)$仅作为总体归一化因子，被高阶twist污染。$\ln z^2$演化可以被精确地定量描述，从而可以从比值函数中提取PDF。

\section{OPE的Wilson系数：从微扰计算到重正子分析}

\subsection{Wilson系数的单圈结构}

准PDF（或spatial correlator）的裸单圈结果在维数正规化（$d=4-2\epsilon$）下具有以下结构\cite{Izubuchi2018}：

\begin{equation}
\begin{aligned}
\tilde{Q}^{(1)}(\zeta, z^2, \epsilon) &= \frac{\alpha_s C_F}{2\pi} e^{\epsilon\gamma_E}\left(\frac{\mu|z|}{2}\right)^{2\epsilon} \Bigg\{ \frac{3}{2}\left(\frac{1}{\epsilon_{\text{UV}}} - \frac{1}{\epsilon_{\text{IR}}}\right) e^{-i\zeta} \\
&\quad + (-1)\frac{\Gamma(2-\epsilon)}{\epsilon_{\text{IR}}} \frac{(1-i\zeta-e^{-i\zeta})}{\zeta^2} + \text{有限项} \Bigg\}。
\end{aligned}
\label{eq:one_loop_bare}
\end{equation}

\textbf{发散结构分析：}
\begin{itemize}
    \item \textbf{$1/\epsilon_{\text{UV}}$项：}由夸克自能和顶点修正的紫外发散产生。在$\overline{\text{MS}}$方案中通过标准抵消项$\delta\tilde{Q}^{(1)} = \frac{\alpha_s C_F}{2\pi} e^{-i\zeta} \frac{3}{2}\frac{1}{\epsilon_{\text{UV}}}$来减除；
    \item \textbf{$1/\epsilon_{\text{IR}}$项：}由共线发散产生。它们在准PDF和光锥PDF之间\textbf{精确匹配}——因此，在匹配核$C = C_{\text{quasi}} - C_{\text{lightcone}}$中彼此抵消；
    \item \textbf{有限项：}包含DGLAP演化对数$\ln(\mu^2 z^2)$和常数项，构成了匹配核的核心。
\end{itemize}

\subsection{重正子与Wilson系数的阶乘发散}

\textbf{这是OPE中与格点精度控制关系最密切的前沿课题。}Braun、Vladimirov和Zhang\cite{Braun2019power}的深入分析揭示了：$\overline{\text{MS}}$方案中的Wilson系数$H(\alpha) = \delta(1-\alpha) + \sum_k h_k a_s^{k+1}$是一个\textbf{渐近级数}——其系数在大阶数下呈阶乘增长：

\begin{equation}
h_k \sim k! \left(\frac{\beta_0}{2\pi}\right)^k \quad (k \to \infty)。
\label{eq:factorial_growth}
\end{equation}

\textbf{重正子的Borel变换分析：}定义Borel变换$\mathcal{B}[H](w) = \sum_k h_k w^k/k!$，分阶乘增长的系数对应于Borel平面上的\textbf{极点奇异点}。

对于准PDF算符，Borel变换的奇异点结构为\cite{Braun2019power}：

\begin{enumerate}
    \item \textbf{紫外重正子（$w = 1/2$）：}来自Wilson线自能中大动量区域的贡献，与线性UV发散一一对应：
    \begin{equation}
    \mathcal{B}[H] \xrightarrow{w \to 1/2} \frac{-4C_F}{w - 1/2} \sqrt{X}, \qquad X = \frac{-v^2 z^2 \mu^2 e^{5/3}}{4}。
    \end{equation}

    \item \textbf{红外重正子（$w = 1, 2, 3, \ldots$）：}领头IR重正子位于$w = 1$：
    \begin{equation}
    \mathcal{B}[H_\parallel](w) \xrightarrow{w \to 1} \frac{-4C_F}{1-w} [\alpha + \bar{\alpha}\ln\bar{\alpha}] X。
    \label{eq:leading_IR_renormalon}
    \end{equation}
    这正是\textbf{twist-3（领先幂次）修正}的微扰信号。
\end{enumerate}

\textbf{重正子模糊性与非微扰修正的匹配：}位于积分路径上的IR重正子极点意味着微扰级数的和是\textbf{病态定义的}（Borel不可求和）。标准的处理方式是将这一模糊性估计为：

\begin{equation}
\delta H(w_0) = -\pi \frac{1}{\beta_0} e^{-w_0/(\beta_0 a_s)} \operatorname*{Res}_{w=w_0} \mathcal{B}[H](w) \propto \left(\frac{\Lambda^2}{\mu^2}\right)^{w_0}。
\label{eq:renormalon_ambiguity}
\end{equation}

对于$w_0 = 1$（领头IR重正子），模糊性$\propto \Lambda^2/\mu^2 \sim \Lambda^2/(P^z)^2$ — \textbf{恰好是twist-3幂次修正的量级}。

\section{从OPE到幂次修正的完整预测}

\subsection{准PDF（qPDF）中的幂次修正}

Braun、Vladimirov和Zhang\cite{Braun2019power}利用重正子分析给出了\textbf{准PDF中幂次修正的完整函数形式}。对于``纵向''投影（$\Gamma = \gamma^t$），功率修正具有以下结构\cite{Braun2019power}：

\begin{definition}[准PDF的领先幂次修正]
\begin{equation}
\boxed{Q_\parallel(x, p) = q(x) \left\{1 + \mathcal{O}\left(\frac{\Lambda^2}{p^2} \cdot \frac{1}{x^2(1-x)}\right)\right\}}。
\label{eq:power_correction_qPDF}
\end{equation}
\end{definition}

\textbf{这一结果的物理含义：}

\begin{itemize}
    \item \textbf{小$x$增强因子$1/x^2$：}幂次修正在小$x$区域被强烈增强——这意味着使用LaMET在有限$P^z$下提取小$x$PDF面临原理性的精度限制；
    \item \textbf{大$x$增强因子$1/(1-x)$：}在$x \to 1$时同样被增强。对于$q(x) \sim (1-x)^\beta$，修正的有效强度被额外放大为$(1-x)^{\beta-1}$；
    \item \textbf{对比靶质量修正：}靶质量修正不显示$1/x^2$增强——这是区别``真实的''高阶twist修正和运动学靶质量修正的重要判据。
\end{itemize}

\textbf{更加具体的函数形式：}从重正子分析导出的积分表示\cite{Braun2019power}：

\begin{equation}
\begin{aligned}
Q_\parallel(x, p) &= q(x) - \kappa \frac{v^2\Lambda^2}{(pv)^2} \left(\frac{d}{dx}\right)^2 \int_{|x|}^1 \frac{dy}{y} (y + \bar{y}\ln\bar{y}) q\left(\frac{x}{y}\right) \\
&= q(x) - \kappa \frac{v^2\Lambda^2}{x^2(pv)^2} \Bigg\{ \int_{|x|}^1 \frac{dy}{y} \frac{y^2}{[1-y]_+} q\left(\frac{x}{y}\right) + q(x) - |x|q'(x) \Bigg\}。
\end{aligned}
\label{eq:qPDF_power_correction_integral}
\end{equation}

其中$\kappa$是$\mathcal{O}(1)$的非微扰参数（重正子分析给出$|\kappa| \approx 1.5$，但确切值需要格点数据的非微扰确定）。

\textbf{零动量归一化对幂次修正的改变：}将qITD归一化到零动量处的单位值（即ratio方案），会\textbf{根本性地改变}幂次修正的结构\cite{Braun2019power}：在$x \lesssim 0.5$时压制修正，但在$x \gtrsim 0.5$--$0.6$时\textbf{强烈增强}。这使得ratio方案不适合用于提取大$x$区域的PDF。

\subsection{赝PDF（pPDF）中的幂次修正}

赝PDF中幂次修正的行为与qPDF有本质不同\cite{Braun2019power}：

\begin{definition}[赝PDF的领先幂次修正]
\begin{equation}
\boxed{\mathcal{P}(x, z) = q(x) \left\{1 + \mathcal{O}\left(z^2\Lambda^2(1-x)\right)\right\}}。
\label{eq:power_correction_pPDF}
\end{equation}
\end{definition}

\textbf{pPDF相对于qPDF的独特优势：}
\begin{itemize}
    \item 在大$x$处，幂次修正是$\mathcal{O}(1-x)$\textbf{压低}的（而非如qPDF中$1/(1-x)$增强）。这意味着\textbf{赝PDF在大$x$区域的系统精度优于准PDF}；
    \item 这一压低性质不被靶质量修正分享（靶质量修正在pPDF中也是$\mathcal{O}(1-x)$压低，但在qPDF中是增强的）；
    \item 然而，零动量归一化会移除这一压低——因此，为了充分利用pPDF的大$x$优势，应避免使用零动量归一化，而采用真空矩阵元归一化。
\end{itemize}

\subsection{靶质量修正：来自运动学幂次展开的贡献}

除了``真实的''高阶twist修正外，OPE还产生\textbf{靶质量修正}（target mass corrections, TMCs）——这些是来自有限核子质量$M^2$的运动学效应。Braun等人\cite{Braun2019power}从光线OPE导出了完整的TMC：

\begin{theorem}[qPDF的靶质量修正]\cite{Braun2019power}
\begin{equation}
\begin{aligned}
Q_\parallel(x, p) &= q(x) + \frac{1}{4}\frac{m^2 v^2}{(pv)^2} [x q'(x) + q(x)] + \mathcal{O}(m^4/p^4), \\
Q_\perp(x, p) &= q(x) + \frac{1}{4}\frac{m^2 v^2}{(pv)^2} [x q'(x) + 3q(x)] \\
&\quad - \frac{1}{2}\frac{m^2 v^2}{(pv)^2} \theta(|x| < 1) \int_{|x|}^1 \frac{dy}{y} q\left(\frac{x}{y}\right)。
\end{aligned}
\label{eq:TMC_qPDF}
\end{equation}
\end{theorem}

\textbf{TMC与``真实''高阶twist的区别：}
\begin{itemize}
    \item TMC在大$x$处也是$\mathcal{O}(1/(1-x))$增强的（对qPDF），但不存在$1/x^2$增强；
    \item TMC仅依赖于核子质量$M^2$，而不依赖于$\Lambda_{\text{QCD}}$——它们可以在不引入额外非微扰参数的情况下\textbf{被精确计算和减除}；
    \item 当$P^z \gtrsim 2$~GeV时，TMC对$M \approx 1$~GeV的相对贡献为$\mathcal{O}(M^2/(P^z)^2) \sim 25\%$——在精度要求下必须被认真对待。
\end{itemize}

\section{Twist展开与幂次修正的级联结构}

\subsection{Twist的严格定义}

在OPE中，一个算符的\textbf{twist}（扭曲度）定义为：

\begin{definition}[Twist]
\begin{equation}
\boxed{\text{twist} = \text{dimension} - \text{spin}},
\label{eq:twist_definition}
\end{equation}
其中dimension是算符的质量量纲（以能量为单位），spin是其Lorentz自旋。
\end{definition}

\textbf{在LaMET展开中的角色：}在$P^z \to \infty$极限下，不同twist的算符对关联函数的贡献按$(1/P^z)^{t-2}$的幂次被压低。因此：

\begin{table}[H]
\centering
\caption{Twist展开的级联结构}
\label{tab:twist_hierarchy}
\begin{tabular}{clcl}
\toprule
Twist & 名称 & 压低因子 & 物理内容 \\
\midrule
2 & 领先twist（leading twist）& $\mathcal{O}(1)$ & 标准PDF（部分子模型） \\
3 & 次领头twist & $\mathcal{O}(1/P^z)$ & 线性幂次修正、重正子模糊性 \\
4 & 次次领头twist & $\mathcal{O}(1/(P^z)^2)$ & 靶质量修正、$z^2\Lambda_{\text{QCD}}^2$项 \\
\vdots & \vdots & \vdots & \vdots \\
\bottomrule
\end{tabular}
\end{table}

\subsection{Twist-3修正与重正子模糊性的精确对应}

\textbf{这是OPE分析的深刻结论}\cite{Zhang2023RGR,Braun2019power}：微扰匹配核中IR重正子的模糊性（$\sim \Lambda^2/(P^z)^2$）和twist-3非微扰矩阵元对关联函数的贡献具有\textbf{完全相同的量级和函数形式}。这意味着：

\begin{quote}
``The twist-two contributions themselves are ambiguous up to higher-twist contributions, and the twist-three parameter $m_0(\tau)$ depends on the resummation/regularization method for the leading renormalon series in $C_k(\alpha_s)$.''\cite{Zhang2023RGR}
\end{quote}

\textbf{这一结论对格点精度控制的关键指导：}仅当微扰匹配核中的重正子发散被一致地指定正规化方案（如主值PV方案），twist-2的匹配和twist-3的修正之间才能实现干净的分离。否则，对两种贡献的任何分离都是方案依赖的。

Zhang等人\cite{Zhang2023RGR}的RGR+LRR方案正是基于这一洞察：在微扰Wilson系数中同时执行\textbf{重整化群重求和}（RGR，控制对数发散）和\textbf{领先重正子重求和}（LRR，控制$1/P^z$幂次发散），将匹配精度从twist-2提升到twist-3（领先幂次）水平。

\section{总结与展望}

算符乘积展开（OPE）是格点QCD中连接微扰短距离物理与非微扰长距离物理的核心理论框架。本文系统梳理了OPE在格点部分子物理中的四重角色及其交叉应用：

\begin{enumerate}
    \item \textbf{标准短距离OPE}提供了PDF的Mellin矩定义，但格点上只能可靠提取前两三个矩\cite{Ji2013Euclidean}。LaMET通过保持算符的非局域形式绕过了算符混合和幂次发散的困难\cite{Ji2017More}。

    \item \textbf{光线OPE（light-ray OPE）}采用$z = \mathcal{O}(\eta), p = \mathcal{O}(\eta^{-1})$的幂次计数，将leading-twist算符的无穷级数重求和为``光射线''结构\cite{Braun2019power}。这是LaMET和赝PDF因子化定理的理论根基。

    \item \textbf{Izubuchi等人的OPE证明}\cite{Izubuchi2018}为LaMET因子化提供了严格的算符级推导，并证明了准PDF和赝PDF方法的等价性——两者仅是同一OPE在不同空间中的实现。

    \item \textbf{重正子分析}\cite{Braun2019power}揭示了微扰Wilson系数的阶乘发散性质，并建立了IR重正子（$w=1$）与twist-3幂次修正之间的精确对应。这一分析给出了准PDF和赝PDF中幂次修正对Bjorken-$x$的完整函数依赖性——qPDF中$\sim 1/[x^2(1-x)]$，pPDF中$\sim (1-x)$。

    \item \textbf{靶质量修正}（TMC）由运动学$M^2/(P^z)^2$效应产生，不引入额外的非微扰参数，可以被精确计算和减除\cite{Braun2019power}。
\end{enumerate}

\textbf{未来方向：}
\begin{itemize}
    \item 将重正子分析（LRR）系统地推广到胶子PDF和螺旋度PDF的匹配核；
    \item 从格点数据中非微扰地提取twist-3矩阵元$m_0(\tau)$，实现``OPE辅助的''精度匹配；
    \item NNLO Wilson系数的计算（对夸克isovector非极化PDF已完成，对胶子和螺旋度仍为NLO）。
\end{itemize}

\begin{thebibliography}{99}

\bibitem{JaffeManohar1990}
R.~L.~Jaffe and A.~Manohar,
``The $g_1$ problem: Deep inelastic electron scattering and the spin of the proton,''
Nucl.\ Phys.\ B \textbf{337}, 509--546 (1990).
\quad\textbf{[来源:} \texttt{docs/The g1 problem: Deep inelastic electron scattering and the spin of the proton.pdf}\textbf{]}
（含$g_1$求和规则的经典OPE分析及Wilson系数的规范不变性论证）

\bibitem{Izubuchi2018}
T.~Izubuchi, X.~Ji, L.~Jin, I.~W.~Stewart, and Y.~Zhao,
``Factorization theorem relating Euclidean and light-cone parton distributions,''
Phys.\ Rev.\ D \textbf{98}, 056004 (2018), arXiv:1801.03917.
\quad\textbf{[来源:} \texttt{docs/Factorization Theorem Relating Euclidean and Light-Cone Parton Distributions.pdf}\textbf{]}
（含LaMET因子化定理的严格OPE证明、约化ITD的短距展开、准/赝PDF等价性）

\bibitem{Braun2019power}
V.~M.~Braun, A.~Vladimirov, and J.-H.~Zhang,
``Power corrections and renormalons in parton quasi-distributions,''
Phys.\ Rev.\ D \textbf{99}, 014013 (2019), arXiv:1810.00048.
\quad\textbf{[来源:} \texttt{docs/Power corrections and renormalons in parton quasi-distributions.pdf}\textbf{]}
（含光线OPE的完整形式、Borel变换与重正子分析、qPDF/pPDF幂次修正的$x$依赖性预测）

\bibitem{Ji2013Euclidean}
X.~Ji,
``Parton physics on a Euclidean lattice,''
Phys.\ Rev.\ Lett.\ \textbf{110}, 262002 (2013), arXiv:1305.1539.
\quad\textbf{[来源:} \texttt{docs/Parton Physics on Euclidean Lattice.pdf}\textbf{]}
（含twist-2局域算符的构造及Mellin矩定义）

\bibitem{Ji2017More}
X.~Ji, J.-H.~Zhang, and Y.~Zhao,
``More on large-momentum effective theory approach to parton physics,''
Nucl.\ Phys.\ B \textbf{924}, 366 (2017), arXiv:1706.07416.
\quad\textbf{[来源:} \texttt{docs/More On Large-Momentum Effective Theory Approach to Parton Physics.pdf}\textbf{]}
（含非局域算符的幂次发散分析和LaMET的有效场论论证）

\bibitem{Ji2014LaMET}
X.~Ji,
``Parton physics from large-momentum effective field theory,''
Sci.\ China Phys.\ Mech.\ Astron.\ \textbf{57}, 1407 (2014), arXiv:1404.6680.
\quad\textbf{[来源:} \texttt{docs/Parton Physics from Large-Momentum Effective Field Theory.pdf}\textbf{]}

\bibitem{Zhang2023RGR}
R.~Zhang, J.~Holligan, X.~Ji, and Y.~Su,
``Leading power accuracy in lattice calculations of parton distributions,''
Phys.\ Lett.\ B \textbf{844}, 138081 (2023), arXiv:2305.05212.
\quad\textbf{[来源:} \texttt{docs/Leading Power Accuracy in Lattice Calculations of Parton Distributions.pdf}\textbf{]}
（含twist-3精度控制、重正子重求和、$m_0(\tau)$参数的非微扰提取）

\bibitem{Su2023LRR}
Y.~Su, J.~Holligan, X.~Ji, F.~Yao, J.-H.~Zhang, and R.~Zhang,
``Power corrections and renormalons in parton quasi-distributions,''
Nucl.\ Phys.\ B \textbf{991}, 116201 (2023), arXiv:2209.01236.
\quad\textbf{[来源:} \texttt{docs/Power corrections and renormalons in parton quasi-distributions.pdf}\textbf{]}

\bibitem{Radyushkin2018one_loop}
A.~Radyushkin,
``One-loop evolution of parton pseudo-distribution functions on the lattice,''
Phys.\ Rev.\ D \textbf{98}, 014019 (2018), arXiv:1801.02427.
\quad\textbf{[来源:} \texttt{docs/One-loop evolution of parton pseudo-distribution functions on the lattice.pdf}\textbf{]}
（含赝PDF的OPE分析与DGLAP演化）

\bibitem{Balitsky2020}
I.~Balitsky, W.~Morris, and A.~Radyushkin,
``Gluon pseudo-distributions at short distances: Forward case,''
Phys.\ Lett.\ B \textbf{808}, 135621 (2020), arXiv:1910.13963.
\quad\textbf{[来源:} \texttt{docs/Short-Distance Structure of Unpolarized Gluon Pseudodistributions.pdf}\textbf{]}
（含胶子算符的OPE/Wilson系数分类）

\bibitem{Chen2025gluon}
C.~Chen \textit{et al.} (CLQCD, Lattice Parton),
``Unpolarized gluon PDF of the nucleon from lattice QCD in the continuum limit,''
(2025), arXiv:2510.26425.
\quad\textbf{[来源:} \texttt{docs/Unpolarized gluon PDF of the nucleon from lattice QCD in the continuum limit.pdf}\textbf{]}

\bibitem{Li2018multiplicative}
Z.-Y.~Li, Y.-Q.~Ma, and J.-W.~Qiu,
``Multiplicative renormalizability of quasi-parton operators,''
Phys.\ Rev.\ Lett.\ \textbf{122}, 062002 (2019), arXiv:1809.01836.
\quad\textbf{[来源:} \texttt{docs/Multiplicative renormalizability of quasi-parton operators.pdf}\textbf{]}

\bibitem{Huo2021self}
Y.-K.~Huo \textit{et al.} (LPC),
``Self-renormalization of quasi-light-front correlators on the lattice,''
Nucl.\ Phys.\ B \textbf{969}, 115443 (2021), arXiv:2103.02965.
\quad\textbf{[来源:} \texttt{docs/Self-Renormalization of Quasi-Light-Front Correlators on the Lattice.pdf}\textbf{]}

\bibitem{Yao2023}
F.~Yao, Y.~Ji, and J.-H.~Zhang,
``Connecting Euclidean to light-cone correlations,''
JHEP \textbf{11}, 021 (2023), arXiv:2212.14415.
\quad\textbf{[来源:} \texttt{docs/Connecting Euclidean to light-cone correlations From flavor nonsinglet in forward kinematics to flavor singlet in non-forward kinematics.pdf}\textbf{]}

\bibitem{NoteGluonPDFs}
``Note of gluon PDFs,'' 内部笔记。
\quad\textbf{[来源:} \texttt{docs/Note of gluon PDFs.pdf}\textbf{]}

\bibitem{Han2025heavy}
X.-Y.~Han \textit{et al.} (LPC),
``Calculation of heavy meson light-cone distribution amplitudes from lattice QCD,''
(2025), arXiv:2410.18654.
\quad\textbf{[来源:} \texttt{docs/Calculation of heavy meson light-cone distribution amplitudes from lattice QCD.pdf}\textbf{]}

\end{thebibliography}

\end{document}
